%! Piecewise Linear Learning Functions — Unit 8, Lesson 8.1 — Slides (Beamer)
\documentclass[aspectratio=169, 11pt]{beamer}
\usepackage{linalg-beamer}   % provides \burgundyheader{} and \sectionlabel[color]{}
\newcommand{\vv}{\mathbf{v}}
\newcommand{\bb}{\mathbf{b}}
\newcommand{\relu}{\mathrm{ReLU}}

\begin{document}

% TITLE SLIDE (CourseName is not defined in beamer, so the name is literal)
{
\setbeamercolor{background canvas}{bg=burgundy}
\begin{frame}[plain]
  \vspace{2.8cm}\hspace{0.55cm}%
  \begin{minipage}{0.9\textwidth}
    {\color{goldacc}\bfseries\small LESSON 8.1}\\[0.5cm]
    {\color{white}\bfseries\LARGE Piecewise Linear Learning Functions}\\[1.0cm]
    {\color{blushmid}\small Linear Algebra~$\cdot$~Unit 8: Learning from Data}
  \end{minipage}
\end{frame}
}

% HOOK
\begin{frame}[t, plain]
  \burgundyheader{One fold is not enough --- so we add folds}
  \sectionlabel{Hook}
  \begin{columns}[T]
    \begin{column}{0.52\textwidth}
      Last lesson: one \textbf{ReLU} $=$ one \emph{fold}; one \textbf{layer} $\relu(A\vv+\bb)$ bends once.\\[6pt]
      A straight line has one slope forever --- it cannot rise \emph{and} fall. To trace bending data we join
      several straight pieces at \textbf{folds}. Two knobs:\\[4pt]
      \textbf{shift} $=$ \emph{where} the fold sits\quad$\bullet$\quad\textbf{weight} $=$ \emph{how steep}.
    \end{column}
    \begin{column}{0.48\textwidth}
      \centering
      \begin{tikzpicture}[scale=0.5]
        \draw[->, charcoal, thin] (-0.4,0)--(4.6,0) node[right]{\scriptsize $x$};
        \draw[->, charcoal, thin] (0,-0.5)--(0,2.8) node[above]{\scriptsize $y$};
        \draw[ultra thick, royalblue] (0,0)--(2,2) node[right]{\scriptsize $\relu(x)$};
        \draw[ultra thick, burgundy] (0,0)--(2,0)--(3.7,1.7);
        \node[burgundy] at (3.5,0.5) {\scriptsize $\relu(x-2)$};
        \fill[burgundy] (2,0) circle (1.6pt);
        \node[charcoal] at (2.2,-0.9) {\scriptsize the shift slides the fold};
      \end{tikzpicture}
    \end{column}
  \end{columns}
  \vspace{2pt}
  \begin{block}{The question of Lesson 8.1}
    How do shifted ReLUs \textbf{add up} into a function that bends to fit any data?
  \end{block}
\end{frame}

% BUILD THE TENT
\begin{frame}[t, plain]
  \burgundyheader{Building a function: the tent}
  \sectionlabel[goldacc]{Skill}
  \[ F(x)=\relu(x)-2\,\relu(x-2)+\relu(x-4). \]
  \begin{columns}[T]
    \begin{column}{0.5\textwidth}
      Turn on only the ReLUs whose fold is left of $x$:
      \[ F(0)=0,\ \ F(2)=2, \]
      \[ F(4)=0,\ \ F(6)=0. \]
      Rises to a peak, comes back down, stays flat --- a \textbf{tent} from three ReLUs.
    \end{column}
    \begin{column}{0.5\textwidth}
      \centering
      \begin{tikzpicture}[scale=0.5]
        \draw[->, charcoal, thin] (-0.4,0)--(5.6,0) node[right]{\scriptsize $x$};
        \draw[->, charcoal, thin] (0,-0.5)--(0,2.9) node[above]{\scriptsize $F(x)$};
        \draw[ultra thick, burgundy] (-0.3,0)--(0,0)--(2,2)--(4,0)--(5.4,0);
        \fill[burgundy] (0,0) circle (1.6pt);
        \fill[burgundy] (2,2) circle (1.6pt);
        \fill[burgundy] (4,0) circle (1.6pt);
        \node[charcoal] at (2.05,2.5) {\scriptsize peak $(2,2)$};
        \node[charcoal] at (0,-0.85) {\scriptsize fold};
        \node[charcoal] at (2,-0.85) {\scriptsize fold};
        \node[charcoal] at (4,-0.85) {\scriptsize fold};
      \end{tikzpicture}
    \end{column}
  \end{columns}
\end{frame}

% SLOPES + COUNTING
\begin{frame}[t, plain]
  \burgundyheader{Reading slopes and counting pieces}
  \sectionlabel{Skill}
  \begin{itemize}
    \setlength{\itemsep}{6pt}
    \item \textbf{Slope on a piece} $=$ sum of the weights of the ReLUs switched on so far. For the tent:
          \[ 0 \ \xrightarrow{+1}\ 1 \ \xrightarrow{-2}\ -1 \ \xrightarrow{+1}\ 0. \]
    \item Each fold \emph{adds} its ReLU's weight to the slope --- up, then down, then flat.
    \item \textbf{Counting rule:} $N$ folds $\Rightarrow N+1$ straight pieces. More ReLUs $\Rightarrow$ more
          folds $\Rightarrow$ a function that bends more.
  \end{itemize}
\end{frame}

% IT IS A LAYER + ROAD AHEAD
\begin{frame}[t, plain]
  \burgundyheader{This is exactly a layer --- and the road ahead}
  \sectionlabel[goldacc]{Payoff}
  \begin{itemize}
    \setlength{\itemsep}{6pt}
    \item A layer $\relu(A\vv+\bb)$ makes these shifted ReLUs: row $i$'s \textbf{bias} places fold~$i$, its
          \textbf{weight} sets the slope there.
    \item \textbf{Fitting data:} choose the weights and biases so the folds and slopes thread the points --- a
          \emph{wider} matrix $A$ has more folds, so it fits \emph{wigglier} data. Shrinking the loss to pick
          them is \textbf{learning} (8.3).
    \item \textbf{Why it matters:} bending to fit data is just shifting and adding ReLUs --- Unit~1 matrices
          plus one fold. \emph{No new machinery.}
  \end{itemize}
  \vspace{2pt}
  \begin{block}{The road ahead}
    \textbf{8.2} convolutional neural nets \;$\bullet$\; \textbf{8.3} minimizing loss by gradient descent
    \;$\bullet$\; \textbf{8.4} mean, variance, covariance.
  \end{block}
\end{frame}

\end{document}
